% 原书物理页=印刷页52--53；3.3 全节.
\section{半单 Lie 代数的分类}\label{sec:ssl-3-3}
设 $\mathfrak g$ 是一个半单 $K$-Lie 代数，$\mathfrak h$ 是一个 Cartan 子代数，$\Phi$ 是 $\mathfrak h$ 在 $\mathfrak g$ 中的根系. 固定 $\Phi$ 的一组基 $\Delta$.
\begin{proposition}\label{prop:ssl-3-3-1}
对每个 $\alpha\in\Delta$，选取直线 $\mathfrak g_{\alpha}$ 的一个生成元 $E_{\alpha}$，并用 $F_{\alpha}$ 表示 $\mathfrak g_{-\alpha}$ 中满足 $[E_{\alpha},F_{\alpha}]=H_{\alpha}$ 的唯一元素. 那么，当 $\alpha$ 取遍 $\Delta$ 时，$\mathfrak g$ 由诸元素 $E_{\alpha}$ 与 $F_{\alpha}$ 生成.
\end{proposition}
\begin{proof}
记 $\mathfrak g'$ 为 $\mathfrak g$ 中由诸 $E_{\alpha}$ 与 $F_{\alpha}$ 生成的 Lie 子代数. 它包含诸 $H_{\alpha}$，而后者构成 $\mathfrak h$ 的一组基，故它包含 $\mathfrak h$. 对正根 $\gamma$，我们对高 $h(\gamma)$ 作归纳来证明 $\mathfrak g_{\gamma}\subset\mathfrak g'$. 若 $h(\gamma)=1$，则 $\gamma$ 是单根，断言由构造即得. 否则，与命题~\ref{prop:ssl-3-2-7} 的证明中一样，存在 $\alpha\in\Delta$ 使 $\langle\alpha,\gamma\rangle>0$，此时 $\beta:=\gamma-\alpha\in\Phi$ 的高为 $h(\beta)=h(\gamma)-1$. 由归纳假设得 $\mathfrak g_{\beta}\subset\mathfrak g'$. 但据命题~\ref{prop:ssl-3-1-5}，有 $\mathfrak g_{\gamma}=[\mathfrak g_{\beta},\mathfrak g_{\alpha}]\subset\mathfrak g'$. 同理可证对负根 $\beta$ 也有 $\mathfrak g_{\beta}\subset\mathfrak g'$.
\end{proof}
\begin{corollary}\label{cor:ssl-3-3-2}
设 $\mathfrak g',\mathfrak h',\Phi',\Delta',E'_{\alpha'}$ 是如上的一组数据. 再设给定了根系的一个同构 $f:\Phi(\mathfrak g,\mathfrak h)\xrightarrow{\ \sim\ }\Phi(\mathfrak g',\mathfrak h')$，使得 $f(\Delta)=\Delta'$. 那么存在唯一的 $K$-Lie 代数同构 $f:\mathfrak g\xrightarrow{\ \sim\ }\mathfrak g'$，使得对每条单根 $\alpha\in\Delta$ 都有 $f(E_{\alpha})=E'_{f(\alpha)}$ 且 $f(H_{\alpha})=H'_{f(\alpha)}$.
\end{corollary}
\begin{proof}
若这样一个同构存在，它还必须把诸 $F_{\alpha}$ 送到诸 $F'_{f(\alpha)}$（因为 $F$ 由等式 $[E,F]=H$ 确定）. 据命题~\ref{prop:ssl-3-3-1}，这就保证了该同构的唯一性.

为证存在性，可以假设 $\Phi$ 不可约，即 $\mathfrak g$ 与 $\mathfrak g'$ 都是单的（习题：请自行说服自己）. 于是考虑 $\mathfrak g\oplus\mathfrak g'$ 中由诸 $E''_{\alpha}:=(E_{\alpha},E'_{f(\alpha)})$ 与诸 $F''_{\alpha}:=(F_{\alpha},F'_{f(\alpha)})$ 生成的子代数 $\mathfrak g''$. 我们将证明投影
\[\mathfrak g''\xrightarrow{\ \pi_1\ }\mathfrak g\qquad\text{与}\qquad\mathfrak g''\xrightarrow{\ \pi_2\ }\mathfrak g'\]
是 $K$-Lie 代数同构. 若确如此，则同构 $\pi_2\circ\pi^{-1}$ 便满足命题的要求.

注意 $\pi_1$ 与 $\pi_2$ 显然是 $K$-Lie 代数的态射，且据前一命题它们是满射. 剩下要证它们是单射. 核 $\operatorname{Ker}(\pi_1)=\mathfrak g''\cap(0\oplus\mathfrak g')$ 是 $0\oplus\mathfrak g'$ 的一个理想（请验证），于是由于 $\mathfrak g'$ 单，它非零当且仅当 $\mathfrak g''$ 包含 $\mathfrak g'=0\oplus\mathfrak g'$，而由 $\mathfrak g''$ 的定义，这等价于 $\mathfrak g''$ 包含 $\mathfrak g\oplus 0$，从而又等价于 $\mathfrak g''=\mathfrak g\oplus\mathfrak g'$. 总之，我们只须证明 $\mathfrak g''\neq\mathfrak g\oplus\mathfrak g'$.

设 $\gamma$ 是 $\Phi$ 的最高根（见上面的习题）. 那么 $\gamma'=f(\gamma)$ 是 $\Phi'$ 的最高根. 选取 $X\in\mathfrak g_{\gamma}\setminus\{0\}$ 与 $X'\in\mathfrak g'_{\gamma'}\setminus\{0\}$，置 $X'':=(X,X')\in\mathfrak g\oplus\mathfrak g'$，并考虑 $\mathfrak g\oplus\mathfrak g'$ 中由所有形如 $\operatorname{ad}F''_{\alpha_m}\operatorname{ad}F''_{\alpha_{m-1}}\cdots\operatorname{ad}F''_{\alpha_1}(X'')$（$m\in\mathbb N$，$\alpha_i\in\Delta$）的元素生成的子 $K$-向量空间 $V$. 由定义可以看出
\begin{enumerate}[label=\roman*)]
\item $V\cap(\mathfrak g_{\gamma}\oplus\mathfrak g'_{\gamma'})=K.X''$，从而 $V\neq\mathfrak g\oplus\mathfrak g'$；
\item $V$ 在诸 $\operatorname{ad}F''_{\alpha}$（$\alpha\in\Delta$）下稳定.
\end{enumerate}
我们来证明 $V$ 在诸 $\operatorname{ad}E''_{\alpha}$ 下也稳定. 为此，注意每当 $\alpha\neq\beta$，$\operatorname{ad}E''_{\alpha}$ 与 $\operatorname{ad}F''_{\beta}$ 交换，因为此时 $[E''_{\alpha},F''_{\beta}]\in\mathfrak g_{\alpha-\beta}\oplus\mathfrak g'_{f(\alpha)-f(\beta)}=0$（由于 $\alpha-\beta$ 不是根）. 于是，利用等式 $\operatorname{ad}E''_{\alpha}\operatorname{ad}F''_{\alpha}=\operatorname{ad}F''_{\alpha}\operatorname{ad}E''_{\alpha}+\operatorname{ad}_{[E''_{\alpha},F''_{\alpha}]}$ 以及 $\operatorname{ad}E''_{\alpha}(X'')=0$（由 $\gamma$ 的极大性，$\gamma+\alpha$ 不是根），一个归纳论证把问题归结为证明 $V$ 在诸 $\operatorname{ad}H''_{\alpha}$ 下稳定，其中 $H''_{\alpha}=[E''_{\alpha},F''_{\alpha}]$. 而这又由同一个归纳论证得出，它基于等式\footnote{译注：下一式交换关系中的正号为原文记号，此处照录；由关系 $(\mathrm R2)$，$[H_\alpha,F_\beta]=-\alpha^\vee(\beta)F_\beta$，相应算子关系的该项应带负号.}
\[\operatorname{ad}H''_{\alpha}\operatorname{ad}F''_{\beta}=\operatorname{ad}F''_{\beta}\operatorname{ad}H''_{\alpha}+\alpha^{\vee}(\beta)\operatorname{ad}F''_{\beta}\]
以及 $X''$ 是 $\operatorname{ad}H''_{\alpha}$ 的特征向量这一事实.

于是 $V$ 在 $\operatorname{ad}(\mathfrak g'')$ 下稳定. 倘若有 $\mathfrak g''=\mathfrak g\oplus\mathfrak g'$，那么 $V$ 就将是 $\mathfrak g\oplus\mathfrak g''$\footnote{译注：原文此处作 $\mathfrak g\oplus\mathfrak g''$；由上下文应为 $\mathfrak g\oplus\mathfrak g'$，因为在假设 $\mathfrak g''=\mathfrak g\oplus\mathfrak g'$ 之下 $V$ 应是 $\mathfrak g\oplus\mathfrak g'$ 的理想.}的一个真理想. 特别地，$V$ 将是半单的，于是
\[V=[V,V]=[V,\mathfrak g]\oplus[V,\mathfrak g']\]
且 $[V,\mathfrak g]=0$ 或 $[V,\mathfrak g']=0$，从而 $V\subset\mathfrak g$ 或 $V\subset\mathfrak g'$，而这是不可能的，因为 $X''\notin\mathfrak g\cup\mathfrak g'$.
\end{proof}
\begin{numpara}[走向分类]\label{np:ssl-3-3-3}
这个推论给出了一条对所有 $K$-单 Lie 代数进行分类的可能策略.
\begin{enumerate}[label=\roman*)]
\item 对所有不可约根系进行分类.
\item 确定其中哪些确实作为某个单 Lie 代数的根系出现.
\end{enumerate}
我们将在后文处理 i). 点 ii) 有一个令人愉快的回答：所有``抽象''根系都出现在单 Lie 代数中（顺便也由此看出，分类并不依赖于那个被设为特征为零代数闭域的域 $K$）. 事实上，我们是幸运的：根系的分类只出现 $4$ 个无穷族，它们对应于一些已经知道的代数（$\mathfrak{sl}_n$、$\mathfrak{so}_{2n+1}$、$\mathfrak{sp}_{2n}$ 与 $\mathfrak{so}_{2n}$），以及有限多个其他的可能，对后者因此可以考虑``手工''构造相应的 Lie 代数.
\end{numpara}
\begin{numpara}[生成元与关系]\label{np:ssl-3-3-4}
然而，前一命题提示了另一条更概念化、也更一致的途径，来构造与一个根系相联系的 Lie 代数：通过生成元与关系. 该命题为我们提供了生成元 $E_{\alpha},F_{\alpha}$（$\alpha\in\Delta$）. 为写出这些生成元之间已知的关系，把诸 $H_{\alpha}$ 也一并引入会更加顺手. 这些关系可写成
\[
\begin{aligned}
&(\mathrm{R}0)\quad [H_{\alpha},H_{\beta}]=0,\\
&(\mathrm{R}1)\quad [E_{\alpha},F_{\beta}]=\delta_{\alpha\beta}H_{\alpha},\\
&(\mathrm{R}2)\quad [H_{\alpha},E_{\beta}]=\alpha^{\vee}(\beta)E_{\beta}\ \text{且}\ [H_{\alpha},F_{\beta}]=-\alpha^{\vee}(\beta)F_{\beta}.
\end{aligned}
\]
于是考虑由生成元 $E_{\alpha},F_{\alpha},H_{\alpha}$（$\alpha\in\Delta$）及关系 $(\mathrm{R}0)$、$(\mathrm{R}1)$ 与 $(\mathrm{R}2)$ 所确定的代数 $\mathfrak g_0(\Phi)$. 在技术上，它是自由 Lie 代数（生成元集为 $\{E_{\alpha},F_{\alpha},H_{\alpha}\}_{\alpha\in\Delta}$）关于由诸元素 $[H_{\alpha},H_{\beta}]$、$[E_{\alpha},F_{\beta}]-\delta_{\alpha\beta}H_{\alpha}$、$[H_{\alpha},E_{\beta}]-\alpha^{\vee}(\beta)E_{\beta}$ 与 $[H_{\alpha},F_{\beta}]+\alpha^{\vee}(\beta)F_{\beta}$ 所生成的理想的商. 可以证明（例如参见 Humphreys 的书，18.2）：诸 $H_{\alpha}$ 在其中生成了一个具有正确维数的 Abel 子代数 $\mathfrak h$，并且 $\mathfrak g_0(\Phi)$ 是它关于 $\mathfrak h$ 的诸权子空间之和. 但一般说来，权的集合是无穷的，因此 $\mathfrak g_0(\Phi)$ 并不是所要找的那个代数. 我们还缺少一些关系，用来迫使商中非零权的集合恰为 $\Phi$.

为找到这样的关系，回忆：对 $\alpha,\beta\in\Delta$ 有 $\beta-\alpha\notin\Phi$. 因而也有
\[s_{\alpha}(\beta-\alpha)=\beta+\bigl(-\alpha^{\vee}(\beta)+1\bigr)\alpha\notin\Phi.\]
由此得到，在一个以 $\Phi$ 为根系的（候选）Lie 代数 $\mathfrak g$ 中，必须成立下列关系（所谓的 Serre 关系）：
\[
(\mathrm S)\quad \bigl(\operatorname{ad}E_{\alpha}\bigr)^{-\alpha^{\vee}(\beta)+1}(\mathfrak g_{\beta})=0\ \text{且}\ \bigl(\operatorname{ad}F_{\alpha}\bigr)^{-\alpha^{\vee}(\beta)+1}(\mathfrak g_{-\beta})=0.
\]
Serre 于是证明了，这些新关系足以构造出所要的 Lie 代数 $\mathfrak g(\Phi)$. 更精确地说：
\end{numpara}
\begin{theorem*}[半单 $K$-Lie 代数的构造]
i) 设 $\mathfrak g$ 是一个半单 $K$-Lie 代数，$\mathfrak h$ 是一个 Cartan 子代数，$\Phi$ 是 $\mathfrak h$ 在 $\mathfrak g$ 中的根系. 固定 $\Phi$ 的一组基 $\Delta$，并回忆 $\alpha^{\vee}(\beta)=\beta(H_{\alpha})$. 那么 $\mathfrak g$ 由诸 $E_{\alpha},H_{\alpha},F_{\alpha}$（$\alpha\in\Delta$）生成，且其仅有的关系为：对 $\alpha,\beta\in\Delta$，\footnote{译注：原文未在下面的关系 $(\mathrm d)$ 中排除 $\alpha=\beta$，此处照录. 该关系须限于 $\alpha\neq\beta$；否则 $1-\alpha^\vee(\alpha)=-1$，所写负次幂没有这里所需的意义.}
\[
\begin{aligned}
&(\mathrm a)\quad [H_{\alpha},H_{\beta}]=0,\\
&(\mathrm b)\quad [E_{\alpha},F_{\beta}]=\delta_{\alpha\beta}H_{\alpha},\\
&(\mathrm c)\quad [H_{\alpha},E_{\beta}]=\alpha^{\vee}(\beta)E_{\beta}\ \text{且}\ [H_{\alpha},F_{\beta}]=-\alpha^{\vee}(\beta)F_{\beta},\\
&(\mathrm d)\quad \bigl(\operatorname{ad}E_{\alpha}\bigr)^{-\alpha^{\vee}(\beta)+1}(\mathfrak g_{\beta})=0\ \text{且}\ \bigl(\operatorname{ad}F_{\alpha}\bigr)^{-\alpha^{\vee}(\beta)+1}(\mathfrak g_{-\beta})=0.
\end{aligned}
\]
ii) 设 $(V,\Phi,\Phi^{\vee})$ 是一个（既约）根系，$\Delta$ 是 $\Phi$ 的一组基. 那么上述展示定义了一个半单 $K$-Lie 代数.

iii) 上述两个过程诱导出双射
\[
\bigl\{\text{半单 }K\text{-Lie 代数}\bigr\}/\text{同构}\;\longleftrightarrow\;\bigl\{\text{（既约）根系}\bigr\}/\text{同构}.
\]
\end{theorem*}
该结果的一个证明可参看 Serre、Carter 或 Humphreys 的书.
\begin{numpara}[自同构]\label{np:ssl-3-3-5}
设 $(\mathfrak g,\mathfrak h,\Phi,\Delta,(E_{\alpha})_{\alpha\in\Delta})$ 如上. 推论~\ref{cor:ssl-3-3-2}（或上述定理）的另一推论是：$\Phi$ 的任何自同构都诱导 $\mathfrak g$ 的一个自同构.
\end{numpara}
\begin{example*}
考虑 $\Phi$ 的自同构 $\theta:\alpha\mapsto-\alpha$. 于是 $\theta(\Delta)=-\Delta$. 置 $\mathfrak g'=\mathfrak g$、$\mathfrak h'=\mathfrak h$、$\Phi'=\Phi$、$\Delta'=-\Delta$，并取 $E'_{-\alpha}:=-F_{\alpha}$. 于是推论~\ref{cor:ssl-3-3-2} 保证存在 $\mathfrak g$ 的唯一的自同构 $\theta$，使得 $E_{\alpha}\mapsto-F_{\alpha}$ 且 $H_{\alpha}\mapsto-H_{\alpha}$. 这个自同构因而也把 $F_{\alpha}$ 送到 $E_{\alpha}$，\footnote{译注：原文这里的 $E_\alpha$ 不带负号，此处照录. 与前面的 $E_\alpha\mapsto-F_\alpha$、$H_\alpha\mapsto-H_\alpha$ 及所给 $X\mapsto-{}^tX$ 相容的像应为 $-E_\alpha$.}从而是 $\mathfrak g$ 的一个对合. 例如在 $\mathfrak g=\mathfrak{sl}_n$ 的情形，有 $\theta(X)=-{}^{t}X$（转置）. 除非 $n=2$，这个自同构都不是内自同构.
\end{example*}
